\section{[\ 附录四\ ]\ 怎样进行作文基础训练}

作文基础训练是指对写字、遣词、造句以及谋篇布局等基本能力的训练。这些能力的训练从小学就开始了。中学阶段，应该继续培养这些方面的能力；并且逐渐加强对分析能力和概括能力、想象能力和逻辑推理能力、观察能力和表达能力等方面的培养。学生只有在各种单项的作文基础训练中掌握了作文基本功，才有可能写出立意新颖、中心突出、结构完整、条理清楚、语言通顺的成篇文章。

为了给成篇的作文训练打下坚实的基础，我们该怎样进行各种单项的作文基础训练呢？

\textbf{1. 要严格进行字词训练。}

就象从平地盖大厦一样，字词训练就是作文训练的基础石。特别是初中同学，应多作这方面的单项训练。常见的训练方式有：

①\ 听写和默写。选择一些范文的片断经常作听写练习；要求写得正确，写得迅速，有意识地训练自己辨别字形、词义的能力，纠正那种大而化之、不注意字的写法和词的用法的坏习惯。对于学过的文言文或者语体文中的精彩片断、名言警句，自觉背诵，并进行默写。这样不仅能对字词有深刻的印象，而且在潜移默化中增强了写作能力。应该注意，不论听写或默写，都要严格要求自己书写工整，并加上恰当的标点符号，以养成一丝不苟的良好习惯。

②\ 经常翻阅“字词手册”。为了积累丰富的词汇，避免作文时没词儿，可以自备“字词手册”，将所读文章中的生字、难字、形近字、同音字、容易读错写错的字词、成语、以及脍炙人口的名言警句，富有表现力的谚语俗话，分门别类记到手册上，一有空子就翻阅记诵。只要坚持不懈，定有成效，写起文章来就能得心应手。美国大作家杰克·伦敦就是用这个“笨办法”掌握了大量的词汇，终于使他在进行创作时左右逢源，各种词汇纷纷涌来，写出一部接一部的世界名著。

③\ 进行字词练习。为了能够正确使用已经积累的字词，可以作一些有典型意义的填字、选词、改正错别字或不当的词语，辨析字词含义等项的练习。这些练习一定要字不离词，词不离句，句不离段，力求在融会贯通上下文意的前提下进行，而不要一字一词地孤零零地去死记硬背，因为那样做是收不到好效果的。

\textbf{2. 要严格进行句子训练。}

同学们首先要学会分析句子的基本结构，然后整理自己作文中出现的病句，分类指出缺少句子成分、成分之间搭配不当、次序颠倒、重复罗嗦、用词不准等语法错误，以及偷换概念、种属混乱、互相矛盾、轻率概括等逻辑错误，并进行改正。我们之所以强调要用同学们自己作文的病句作辨析改正的对象，是因为这样做了，同学们会感到学的不是死知识，而是活生生的东西，改起来有兴趣，印象深刻，效果好。

在运用已有知识纠正自己作文中病句的同时，还应该进行造句练习，以便把获得的知识变为技能。造句的方式可以多种多样。如用单词造句，联缀两个以上的词造句，模仿课文中的某些句式造句，按规定的句式造句，缩写长句，扩写短句，给一个上句来对一个下句（对对子），变换句子形式，等等。其中以进行变换句式的训练最有实效。据说，俄国大文豪列夫·托尔斯泰在青少年时代就曾在这方面下过苦功，常常把一句话的意思分别用陈述句式、疑问句式、祈使句式、感叹句式或者主动句式、被动句式、倒置句式、把字句式等等多种的句式来表达，因而他在后来的文艺创作中能够得心应手，表意准确、严密、富于变化而又无比生动。

\textbf{3. 要严格进行文意连贯的训练。}

同学们作文中常见的毛病是语无伦次，前言不搭后语。这就是说，单会造句子，而不会逻辑严密地组合句子，仍不能写出一段文意连贯、条理清楚的文字，当然就更难成篇了。这是在进行作文基础训练时必须解决的一个课题。

写文章层次不清，语言混乱，是思维不清的反映。为了学会按照一定的规律思考问题，学会把自己的思想确切地、有条有理地表达出来，我们可以从这样几个方面入手去做：

①\ 范文引导。自习时，有意识地将某一文理严谨的段落中各句之间的关系弄清楚，或将全文的段落层次弄清楚。一定要从反面搞懂为什么不能颠倒的道理，这对训练思维能力颇有实效。对于不同文体的文章，应该抓住它们不同的特点，多想几个为什么，彻底弄清作者的思路。

②\ 加强背诵。学一些逻辑思维的知识固然必要，但更为有效的方法倒是熟读和背诵名作。当代著名文学家巴金，在青年时代能将《古文观止》中的两百多篇佳作背诵出来，这使他终生受用不尽。他说：“现在有两百多篇文章储蓄在我的脑子里面了……这么多具体东西至少可以使我明白所谓‘文章’究竟是怎么一回事，它是有条有理，顺着我们的思路走下来的。”有一个学生的作文，经常语无伦次，经过一年多背诵名篇的锻炼，思路开始明晰了。看来背诵这个办法似乎有点笨，但是只要对名著的内容有所理解了，作文水平就会相应地有所提高。所谓“读书破万卷，下笔如有神”，是很有道理的。

③\ 整理语序。为了训练逻辑思维能力，可以经常进行整理语序的练习。比如将一段结构严谨、层次分明的话打乱句序，或将一篇逻辑严密的文章的段落颠倒颠倒，自己去找出句与句、段与段的内在联系，然后把文章的脉络理顺，并讲出道理来。这样的训练，可以有助于养成动笔之前先动一番脑筋想一想文章脉络的习惯，有助于提高根据材料的内在联系组织文章的能力。

④\ 由点及面。在有了联句成段的能力基础上，可以由小到大，由点到面地联段成篇。这样的训练，对于学习运用概念进行判断和推理，学习分析和综合，是很有帮助的。

\textbf{4. 要严格进行谋篇布局的训练。}
①\ 审题训练。目前作文训练大致有三种类型的命题方法，一是命题作文；二是根据材料命题作文；三是限制范围自己拟题。除了第三种需要根据所定的范围和确立的中心自拟恰当而精炼的题目以外，前两种都有个审题问题。从“十年动乱”前高考的命题作文和“十年动乱”后高考的根据材料命题作文来看，是否能正确审题是一篇作文成败的关键。这就要在平时作好严格的审题训练。

审题训练可采取分析法和比较法进行。分析法是说分析题目本身的结构和含义，比如从语法结构上分析，从关键性的词语来分析，从限制性的修饰词来分析，从词的内在寓意来分析，等等，以确定文章体裁、选材范围、中心思想、写作重点和写作要求。比较法是把一组结构类似或不同、含义相近或相异的题目，拿来进行审察、分析、研究，然后加以归纳，明确它们之间有什么不同点，哪怕是极其细微的差别，由此来确定文章的不同写法。只要经常注意审题训练，使得写作路子对了，就不愁在作文时再出现下笔千言、离题万里、东扯西拉、没有中心的现象了。

②\ 立意训练。立意就是确定文章的中心思想，这是组织材料或选用论据的依据。立意不明或立意不高，正是初学写作者容易犯的毛病。立意不明，将使文章成为材料的堆砌，现象的罗列，记流水帐，立意不高，将使文章停留在就事写事、就事论事上面。这两种现象，都使文章缺少思想价值，缺乏感人力量。

为了使作文能做到立意突出、深刻、新颖，可以进行确定主题的训练。比如请老师提供一篇记叙文的事件材料，自己着手分析这些材料能说明什么问题，有什么思想意义，从而明确其中心思想，并以画龙点睛的写法将文章的尾段补述出来，或者请老师提供一篇议论文的论据材料，自己来归纳其中心论点，并且旗帜鲜明地在首段把它补述出来。

另外，还可以进行编写作文提纲的训练。要求自己在确定文章中心思想以后，写出比较详细、包括各段所用材料的提纲，因为只有粗纲，而无细目，成文之后，仍然可能有游离于立意之外的材料或与立意扣得不紧的闲笔。

③\ 组织材料的训练。写文章应该根据主题的需要，把选来的那些零碎的材料，进行合理的组织安排，构成有机的整体。要求做到：层次清楚，线索分明，详略得当，中心突出，过渡自然，前后照应，工于开端，妙于结尾，使全文结构完整、谨严、有逻辑性。

为了使自己具有这种组织材料的能力，可以进行如下的练习：

填写首段。要求根据供给的材料，添加恰当的开头，使文章结构完整。

填写尾段。要求根据供给的材料，添加恰当的结尾，使文章结构完整。

写结构分析提纲。要求用少量文字把课文的基本内容（主要情节或论述要点）表述出来。这对于加深理解课文内容，特别是理解课文结构，对于培养概括事物的能力，很有帮助。

写结构仿写提纲。要求借鉴课文的结构提纲，自己着手写出一组作文题目相近或类似的写作提纲。写出的提纲要有怎样开头、发展、过渡、结尾的细目，看看是否符合文章的起、承、转、合的基本程式。

改正首尾脱节的毛病。请老师提供首尾不相呼应的材料，自己进行改正。

调整段落次序。请老师提供段落次序颠倒的材料，自己进行改正。记叙文要按空间、时间或人物活动的变化来调整不合逻辑的顺序；议论文要按中心论点、分论点或论据的性质来调整不合逻辑的顺序。

整理零乱材料，加工成文。请老师提供零碎杂乱的有关材料，并进行命题或提出整理成文的要求。自己思考哪些材料紧扣题意，能够说明中心，就集中、概括起来，而把与中心无关的材料舍弃；然后分清主次，理清头绪，按文章的逻辑顺序加以组织安排，做到纲举目张，条分缕析，结构谨严，层次清楚；再加上必要的开头、过渡或结尾，从而组成一篇完整的有条有理的文章。

当然，为了能将学到的语文知识化为技能，还必须把写和听、说、读结合起来进行训练才好。比如复述故事或课文，口头作文，写阅读笔记，作文章摘要，记日记，写作文后记等等，都有利于训练写作基本功，我们可以根据自己的实际情况有意识地去做。
\newpage